\chapter*{Tóm tắt}
\thispagestyle{empty}

Ở thời đại số hiện nay, công nghệ đã và đang trở thành một phần không thể thiếu trong giáo dục, góp phần đổi mới phương pháp dạy và học theo hướng hiện đại, chủ động và hiệu quả hơn. Việc ứng dụng công nghệ không chỉ giúp người học mở rộng kiến thức, tiết kiệm thời gian mà còn tạo điều kiện để tiếp cận bài học một cách trực quan và sinh động. Đối với môn Vật lý 11, các nội dung như dao động điều hòa, sóng cơ, điện trường hay dòng điện không đổi tuy xuất phát từ những hiện tượng quen thuộc trong đời sống hằng ngày, nhưng lại mang tính trừu tượng cao, khiến học sinh khó hình dung nếu chỉ tiếp cận qua sách giáo khoa và lời giảng lý thuyết. Trong bối cảnh đó, các phần mềm mô phỏng và công cụ học tập số có thể đóng vai trò như “cầu nối” giữa lý thuyết và thực tiễn, giúp học sinh quan sát quá trình diễn ra của hiện tượng, thử nghiệm các yếu tố ảnh hưởng và từ đó hiểu bản chất vấn đề một cách sâu sắc hơn. Tuy vậy, việc áp dụng những công cụ số trong học tập đòi hỏi sự chủ động và yêu cầu người dùng có nền tảng về công nghệ, điều này khiến chúng vẫn còn xa lạ và khó tiếp cận đối với phần đông học sinh và giáo viên. 

Xác định được thực tế trên, nhóm thực hiện đề tài "\textbf{CHƯƠNG TRÌNH HỖ TRỢ DẠY VÀ HỌC MÔN VẬT LÝ LỚP 11}" với mục tiêu tạo ra một nền tảng học tập không chỉ mô phỏng các hiện tượng vật lý mà còn hệ thống hóa kiến thức theo chuẩn chương trình sách giáo khoa, từ đó có thể được triển khai chính thức và đồng bộ trong giảng dạy, giúp học sinh và giáo viên hình thành thói quen sử dụng các công cụ số trong quá trình tiếp thu và truyền đạt kiến thức, đồng thời phát huy tối đa các lợi ích của công nghệ trong giáo dục.

\vspace{0.5em}

Sản phẩm phần mềm khi đưa vào vận hành sẽ cần phải đạt được 4 yêu cầu cơ bản sau: (1) Giả lập mô phỏng được xây dựng có thể tương tác, cho phép người dùng thay đổi các tham số vật lý và quan sát kết quả một cách trực quan; (2) Giao diện trình bày nội dung lý thuyết gần gũi dễ sử dụng, bám sát cách tổ chức kiến thức trong sách giáo khoa (SGK) \textit{Chân trời sáng tạo} (3) Hệ thống bài tập đa dạng các cấp độ, những bài tập khó được chọn lọc từ những đề chuyên, giúp học sinh rèn luyện và nâng cao kỹ năng qua các câu hỏi từ cơ bản đến nâng cao, bao gồm bài tập trắc nghiệm, bài tập tính toán và bài tập tình huống thực tế. (4) Kiến trúc phần mềm linh hoạt, đảm bảo tính dễ bảo trì và mở rộng trong tương lai, giúp phần mềm có thể được cập nhật và bổ sung các tính năng mới hoặc nội dung giáo dục mà không làm ảnh hưởng đến hiệu quả sử dụng lâu dài.

\clearpage
\pagenumbering{arabic}
